\section{DiagPriv Transformation Framework}
\label{sec:framework}
The implementation keeps benchmark construction, transformations, prompt
rendering, provider invocation, and oracle scoring separate. A generator
produces an \texttt{Incident} with model-visible records and oracle-only
labels, privacy annotations, expected evidence, and invariant definitions.
Transformations begin from \texttt{Incident.model\_view()} and return a new
\texttt{ModelView}; neither an oracle answer nor the HMAC key is serialized
into the model request. A versioned condition and hash identify each
representation. The fixed baselines below isolate different preservation
choices rather than claiming to be interchangeable security mechanisms.

\paragraph{Redaction and deterministic pseudonyms.}
\texttt{redaction-v1} replaces annotated values with field-class
placeholders. It hides raw literal identity but can collapse distinct
entities and destroy meaningful joins. \texttt{hmac-v1} uses keyed,
namespace-separated HMAC-SHA-256 identifiers; it preserves deterministic
equality within declared namespaces, including selected cross-source joins,
while retaining frequency and linkability. Its operational key is external
to prompts and artifacts. The privacy-exposure analysis uses a separate
public synthetic test-only key and makes no operational key-compromise claim.
\texttt{tokenization-v1} assigns short, deterministic, ephemeral tokens from
a mapping over the complete view. Its mapping is reproducible for the same
input population but is neither a persistent production vault nor
guaranteed globally stable across independently transformed candidate sets.
Both methods can expose candidate-set or frequency information even when raw
values disappear.

\subsubsection*{Shape and time}
\texttt{format-preserving-v1} is deterministic format/constraint-preserving
pseudonymization: it retains measured length, separator positions, and
character classes while changing literal values. It is \emph{not}
cryptographic format-preserving encryption. \texttt{relative-time-v1}
converts annotated timestamps to incident-local millisecond offsets,
removing the original absolute representation but keeping equality,
ordering, gaps, durations, and overlap. Exact offsets could reconstruct
absolute times given an external anchor. \texttt{temporal-ablation-v1}
replaces annotated timestamps with non-temporal markers and removes log-order
leakage; it is a destructive causal control, not an optimal privacy method.
In Incident 001 and 002 the scalar methods leave \emph{unannotated}
timestamps unchanged, a limit measured explicitly below.

These conditions are versioned; no transformation was retuned after observing
the results reported here. The broader DiagPriv ambition is to choose a
least-exposing condition that preserves the declared essential relations for
the task. This draft evaluates the premise with fixed baselines. It does not
present an implemented task-aware planner or an automatic guarantee of
minimal disclosure.
